\documentclass[10pt,a4paper]{amsart}
\usepackage[margin=19mm]{geometry}
\usepackage{amsmath,amssymb,mathtools}
\usepackage[scr=rsfs,cal=euler]{mathalfa}
\usepackage{xfrac}
\usepackage[unicode,hidelinks,bookmarks=false]{hyperref}
\newcommand{\ie}{i.e.\ }
\newcommand{\cf}{cf.\ }
\newcommand{\resp}{respectively\ }
\newcommand{\Subsection}{\subsection}
\providecommand{\nobreakdash}{\nobreak\hbox{-}}
\setlength{\emergencystretch}{2em}
\multlinegap0pt
\usepackage{tcolorbox}
\usepackage{xcolor}
\newtheorem{proposition}{Proposition}
\newcommand{\PP}{\mathbb{P}}
\newcommand{\T}{\mathbb{T}}
\newcommand{\bP}{\textbf{P}}
\title{Cutoff for the mixing time of the Facilitated Exclusion Process}
\author{Brune Massouli\'e}
\date{Source: arXiv:2412.04032v2 (CC BY 4.0)}
\begin{document}
\maketitle

\noindent\emph{Source and licence.} Cutoff for the mixing time of the Facilitated Exclusion Process. Version arXiv:2412.04032v2 (CC BY 4.0), distributed by arXiv under the Creative Commons Attribution 4.0 International licence, http://creativecommons.org/licenses/by/4.0/, as recorded in the arXiv licence metadata for this exact version and shown on the abstract page https://arxiv.org/abs/2412.04032v2. Full licence notice: \texttt{licenses/arxiv-2412.04032-CC-BY-4.0.txt}.\par\smallskip

\noindent\textit{Editorial scope.} Counted unit: the article's proposition prop:link\_fep\_ssep (source0.tex lines 683--688). The article's complete original proof of it is delivered with it (source0.tex lines 689--708, one \texttt{proof} environment of 20 lines). No mathematics is written, repaired or re-derived: the statement and the proof are the article's own lines, in the article's own notation, and no retained line is substituted or re-flowed.  No further source-local context is needed: every cross-reference of every retained line resolves inside the two delivered spans.  Nothing was omitted from any retained span, so the package carries no omission record.  No retained line contains a citation, so the wrapper carries no bibliography and no bibliography entry.  No retained line contains a figure environment, an image-inclusion command or drawing code, so no image is delivered and the package declares no asset dependency.  Editorial formatting only, in addition: an AMS \texttt{amsart} preamble that restates the article's own declarations and macros used by the retained lines, namely the environments \texttt{proposition} on separate counters, together with the standard packages those lines need.  The retained environments print in this document as Proposition 1 in order of appearance, whereas the article prints the same items with the numbering of its own full document.  The complete native source of the article as distributed in the arXiv e-print for this exact version is bundled unchanged as \texttt{sources/169651/source0.tex}.
\begin{proposition}[Link between particle density in the SSEP and the FEP] \label{prop:link_fep_ssep}
For all $s, P \ge 1$,
\begin{multline}
\PP_\eta\left(\exists \hbox{ a segment } I \subset \T_N,  \left||\eta_{|I}(t_1) | - \frac{K}{N} |I|\right| \ge s \frac{K}{N} \sqrt{P}\right) \\ \le \textnormal{\bP}_\sigma\left(\exists k,l, \left|\sum_{j=k}^{l} \left(\sigma_j(t_1) - \frac{P}{K} \right) \right|\ge s\sqrt{P} -1\right).
\end{multline}
\end{proposition}

\begin{proof}
Notice that a segment $J$ of size $j$ in $\sigma(t_1)$ corresponds to $j$ consecutive particles in $\eta(t_1)$, which are contained in a segment of size $\sum_{k\in J} 2 - \sigma_k(t_1)$. Indeed, if there is a particle at one site in $\sigma(t_1)$, it means there is no space between the corresponding particle in $\eta(t_1)$ and its right neighbour, so this corresponding particle occupies a space of 1 in $\eta(t_1)$. If there is no particle at a site of $\sigma(t_1)$, it means the corresponding particle in $\eta(t_1)$ is followed by an empty site in $\eta(t_1)$, so it occupies 2 spaces in $\eta(t_1)$. 

So given a segment $J$ in the SSEP, there is a corresponding segment in the FEP, starting with a particle, of length $\sum_{k\in J} \left(2-\sigma_j(t_1)  \right)$. Now, replacing $P$ by $2K-N$, notice that:

\begin{multline}\bP_\sigma\left(\exists k, l : \left|\sum_{j=k}^{l} \left(\sigma_j(t_1) - \frac{2K-N}{K} \right) \right|\ge s\sqrt{P} \right) \\ = \bP_\sigma\left(\exists k, l, \left| \frac{N}{K}(l-k +1) - \sum_{j=k}^{l} \left(2 - \sigma_j(t_1)\right)\right| \ge s\sqrt{P}\right),
\end{multline}
which is  therefore the probability that in the FEP, there exists a segment $I$ starting with a particle such that $\left| |\eta_{|I}(t_1)| - \frac{K}{N} |I| \right| \ge s\frac{K}{N} \sqrt{P}$. 

If there exists a segment $I = [x,y]$ starting with an empty site
such that $\left| |\eta_{|I}(t_1)| - \frac{K}{N} |I| \right| \ge s \frac{K}{N}\sqrt{P}$, then by ergodicity, the segment $I' = [x+1, y]$ starts with a particle, and we have $\left| |\eta_{|I'}(t_1)| - \frac{K}{N} |I'| \right| \ge \frac{K}{N} (s\sqrt{P}-1)$. So, 
\begin{align}
&\,\PP_\eta\left(\exists \hbox{ a segment } I, \left| |\eta_{|I}(t_1)| - \frac{K}{N} |I| \right| \ge s\frac{K}{N} \sqrt{P}\right)
 \nonumber\\
\le &\,\PP_\eta\bigg(\exists \hbox{ a segment } I \hbox{ starting with a particle }, \nonumber\\
& \qquad\qquad\qquad\qquad\qquad\qquad \left| |\eta_{|I}(t_1)| - \frac{K}{N} |I| \right| \ge \frac{K}{N}(s\sqrt{P}-1)\bigg) \nonumber\\
=& \,\bP_\sigma\left(\exists \hbox{ a segment }J \subset \T_K, \left|\sum_{k\in J}\left(\sigma_k(t_1) - \frac{P}{K}\right)\right| \ge s\sqrt{P}-1\right).
 \end{align}

\end{proof}

\end{document}
